\section{Leakage Audit}
\label{sec:audit}

\subsection{The original split fails}

The first audit pass over the pooled corpus returned a clear failure. Of the
evaluation images, $63.8\%$ had a training image within perceptual-hash Hamming
distance $2$, and $76$ Roboflow base names appeared in more than one split
(train$\cap$valid $37$, train$\cap$test $31$, valid$\cap$test $10$). A montage
of the base-name collisions shows the same physical scene, the same mine
placement, and different export bytes: these are re-exports, not distinct
photographs. The mechanism is augmentation before splitting, present in the
upstream generation scripts. The test split also contained zero negative images,
so false-alarm rate was undefined. We therefore discarded the split entirely
rather than attempting to patch it.

\subsection{Operational leak definition}

A validation or test image $x$ is a leak if any training image shares any of:

\begin{enumerate}[leftmargin=1.4em,itemsep=1pt]
\item the exact MD5 digest;
\item the Roboflow base key (filename with the \texttt{*.rf.<md5>.jpg} export
suffix stripped);
\item the same source id \emph{and} a 64-bit perceptual hash \citep{zauner2010phash}
within Hamming distance $2$ under all eight dihedral transforms, \emph{and} a
$32\times32$ RGB mean-squared error below $900$.
\end{enumerate}

Condition 3 is a conjunction for a reason. Perceptual hash alone over-fires on
repetitive terrain: grass and dirt produce Hamming collisions that are not the
same scene. The MSE term separates the two populations, and manual review sets
the cut. Every Hamming-$\le 2$ candidate---$146$ pairs across two rounds---was
rendered side by side and inspected. Confirmed same-scene pairs occurred at
$\mathrm{MSE} \le 817$; terrain-texture coincidences at
$\mathrm{MSE} \ge 958.6$. The gap is narrow but unambiguous on montage, and the
cut at $900$ sits inside it.

\subsection{Rebuilding and freezing}

Groups were formed from base key, capture session, and same-source hash edges,
yielding $1{,}048$ groups, then assigned whole to a split. Two filtering passes
removed $30$ evaluation images: $13$ at $\mathrm{MSE}\le 100$, then $17$ more at
$\mathrm{MSE}\le 900$, including a cross-folder duplicate
(\texttt{ds0\_0043}$\leftrightarrow$\texttt{ds1\_0043}). The split is
$19$ validation and $11$ test images lighter as a result. The residual $99$ of
$370$ evaluation images with a Hamming-$\le 2$ training neighbour all have
$\mathrm{MSE} \ge 958.6$ and were verified as texture false positives. True leak
rate: $0/370$.

\begin{table}[t]
\centering
\small
\caption{Leakage audit outcome for the frozen split. ``Residual candidates''
are Hamming-$\le2$ pairs rejected by the pixel-MSE test and manual review.}
\label{tab:audit}
\begin{tabular}{lr}
\toprule
quantity & value \\
\midrule
Images (train / valid / test) & 3{,}907 / 181 / 189 \\
Instances & 11{,}550 / 419 / 253 \\
Negatives & 20 / 7 / 5 \\
Exact cross-split duplicate groups & 0 \\
Base-key collisions (train-valid / train-test / valid-test) & 0 / 0 / 0 \\
Residual near-duplicate candidates & 99 (all MSE $\ge 958.6$) \\
Evaluation images removed as leaks & 30 (19 valid, 11 test) \\
Confirmed true leak rate & $0/370 = 0.0\%$ \\
Manifest SHA-256 & \texttt{748322d6\ldots c80de} \\
\bottomrule
\end{tabular}
\end{table}

\subsection{What the leak was worth}
\label{sec:leakvalue}

An audit is only useful if it changes something. We therefore trained a YOLOv8s
detector on the \emph{old} split (the ``leaky'' model) and compared it with the
clean-trained seed-0 detector on the clean test set. This is the natural
before/after experiment, and it has a trap: $140$ of $189$ clean test images
($74.1\%$) share a v2 group with an image the leaky model trained on. Scoring the
leaky model on the full test set therefore measures memorisation.

The results are in Table~\ref{tab:leaky}. On its own split the leaky model
reaches $0.970$ mAP$_{50}$, but that split has zero negatives, so false-positive
rate is undefined and the number cannot be compared with anything in this paper.
On the full clean test set it reaches $0.994$, comfortably above the clean
model's $0.974$. On the $49$ contamination-free images the two are tied:
$0.9950$ against $0.9947$ mAP$_{50}$, and $0.725$ against $0.701$ mAP$_{50\text{--}95}$.
The $49$-image subset is small and the comparison is therefore low-powered, but
the direction is unambiguous: the apparent $+0.020$ advantage is contamination,
and the honest estimate of the leak's effect on this corpus is close to zero.

\begin{table}[t]
\centering
\small
\caption{Leaky-versus-clean comparison. The leaky model's operating confidence
($0.065$) was chosen on the old validation split; the clean model's ($0.02$) on
the new one. P and FP/image are therefore not threshold-matched and mAP is the
fair column.}
\label{tab:leaky}
\begin{tabular}{llcccc}
\toprule
model & evaluation set & mAP$_{50}$ & mAP$_{50\text{--}95}$ & P & FP/img \\
\midrule
leaky v1 & v1 test (own split) & 0.9697 & 0.6436 & 0.834 & n/a (0 neg) \\
leaky v1 & v2 test, full (74.1\% contaminated) & 0.9940 & --- & 0.857 & 0.222 \\
clean v2 s0 & v2 test, full & 0.9741 & 0.6924 & 0.427 & 1.751 \\
\midrule
leaky v1 & v2 test, contamination-free ($n=49$) & 0.9950 & 0.7250 & 0.863 & 0.204 \\
clean v2 s0 & v2 test, contamination-free ($n=49$) & 0.9947 & 0.7010 & 0.354 & 2.347 \\
\bottomrule
\end{tabular}
\end{table}

This is a negative result and we treat it as one. The literature on leakage
\citep{kapoor2023leakage,recht2019imagenet} prepares one to expect a large
inflation. Here the clean and leaky models converge on the leak-free images, which
suggests that for this task the leaked material mostly taught the model the same
thing the clean data does. That is not a reason to skip the audit---the old split
remains unusable, because its selection split was contaminated and its test split
had no negatives---but it is a reason to report the size of the effect honestly
rather than using the audit as a rhetorical amplifier. It also has a practical
implication: an audit that only reports ``we fixed leakage'' without measuring the
contamination-free delta leaves the reader unable to judge whether the fix
mattered.
